\chapter{Evaluation Plan}
\label{sec:evaluation}

The design's soundness can be argued at a desk; its value cannot. Whether the rule is worth adopting
depends on how often it rejects code that people actually write, how many of those rejections are real
sharing rather than limits of the checker, and whether the widened demand set is acceptable. We know of
no published false-error rate for an inferred uniqueness checker in a functional language. This chapter
describes the experiment we propose to run before deciding. \emph{It has not been run, and we report no
results.}

The experiment follows the path by which beni's write-set analysis was introduced
(\S\ref{sec:bg-browser}): build the analysis, run it as a report over real programs, classify every
rejection, and decide on the numbers. Following Wrenn and Krishnamurthi's view of errors as classifiers
\citep[\S2]{wrenn2017classifiers}, the experiment measures the precision of the classifier the checker
would be.

\section{Report mode}
\label{sec:eval-report}

The analysis is built first as a pass that only reports. It runs during checking, prints a dump---for
instance under \code{beni dump --stage=ownership}---and changes no emitted byte, so no test program's
output or recorded hash moves; nothing reads the dump but people and tests. The dump prints, per module:
\begin{Verbatim}
summary <Decl>
  param <i>.<path> <borrowed|shared|consumed> [reason <tag>]
  result <path> {<param paths>|fresh}
  fun <j> calls <0|1|many> escapes <no|yes>
demand <Decl>:<inst> <primitive> operand <cells>
  unique                                              -- accepted
  shared <holder> kept-at <inst> used-at <inst> exact  -- real sharing
  shared <holder> … approx <limit-tag>                 -- possible sharing (a limit)
  unknown <limit-tag>                                  -- Known or Unescaped failed
copy <Decl>:<inst> of-unique                           -- a copy of a list already unique
prepend <Decl>:<inst> <in-building-loop|loose>
cap <which> fired at <Decl>
\end{Verbatim}
The dump is deterministic (\S\ref{sec:inf-modular}), is compared byte for byte by the determinism test,
and is itself tested against expected output for a set of small programs, one per transfer rule and per
side condition, each written to fail before the rule exists.

\section{Two demand sets}
\label{sec:eval-demands}

A hidden flag selects the demand set, so that one run reports both the rule as stated (the nine writers
of \S\ref{sec:model-demands}) and the widened rule (change~1 of Chapter~\ref{sec:disc-options}: plus
\code{map}, \code{indexedMap}, \code{filter}, \code{filterMap}, \code{reverse}, \code{sort}, \code{sortBy},
\code{sortWith}, \code{take} and \code{slice}). Every count is reported for both.

\section{Corpora and counts}
\label{sec:eval-counts}

The corpora are beni's core library; the compiler's test corpus (programs that run under Node,
programs that run in a page, and the write-set analysis's own tests); the four TodoMVC builds; the
Conduit application; and a table-rendering benchmark application with its parameter sweeps. TEA programs
are analysed under the hand-over obligations of \S\ref{sec:hard-tea}, with the entry summary of the
model computed by the entry fixpoint. For each corpus and in total we report:

\begin{center}\small
\begin{tabularx}{\linewidth}{@{}>{\raggedright\arraybackslash}p{0.27\linewidth}X@{}}
\toprule
Count & Meaning \\
\midrule
demands & update sites, per demand set \\
accepted & unique at the site \\
rejected, exact & real sharing, by holder kind: a later local read, a closure, a command or fiber, a
container, a payload, the view \\
rejected, approximate & by limit tag: \code{foreign}/\code{Js}, \code{Ref}, path cut, join, pessimistic
summary, $\mathit{calls} = \mathit{many}$, elements may share \\
exclusive containers & container paths that carry $\excl$ at a destructive read, against those that do
not, and what cleared the flag \\
copies of unique lists & \code{List.copy} calls the checker would not have needed \\
loose prepends & \code{[ x, …xs ]} sites outside a building loop (the $O(n)$ ones) \\
summaries changed per edit & under the fixed set of edit classes already used to measure interface
churn, applied to every public declaration of the core library, how many interfaces change \\
check time & the pass's time beside checking, cold, optimised build, on a benchmark corpus and on
Conduit, from the compiler's self-profiling events \\
entry fixpoint time & the whole-program entry computation's time, beside the write-set analysis's \\
widest $w$ & the largest number of tracked paths in any declaration \\
\bottomrule
\end{tabularx}
\end{center}

\section{Hand classification}
\label{sec:eval-hand}

Every exact rejection is read by a person and classified as \emph{intended sharing} (the program keeps
a version on purpose: an undo, a snapshot), \emph{accidental sharing} (the program would have been wrong
had the write been in place---the bug class the rule exists to catch), or \emph{avoidable} (a reordering
or a fold would make the operand unique). Every approximate rejection is classified by whether a known
extension---a discipline for mutable cells, a deeper path cut, a path-sensitive join, a summary on a
platform \code{foreign}---would accept it. The model is the hand analysis in Cann, Feo and DeBoni's report
on the Sisal compiler, which tabulates copy operations in five application programs before and after
optimisation---293 before, none unconditional and 29 conditional after---and explains the remaining
conditional copies in its prose by row sharing \citep{cann1990sisal}.

\section{Decision criteria}
\label{sec:eval-decisions}

The numbers inform three decisions, in this order. First, the \emph{demand set}: under the rule as
stated the two applications are untouched (\S\ref{sec:disc-apps}), and under the widened rule every list
edit in them is a demand, so the rejection rate under the widened set decides whether widening is
acceptable. Second, \emph{error or warning}, on the split between exact and approximate rejections: the
soundness argument holds for both, and the choice is about the cost model and the principle of
Chapter~\ref{sec:discussion}. Third, which precision extensions to build, on the frequency of each
approximate tag---a path-sensitive join only if the join tag is frequent, a discipline for mutable
cells only if the \code{Ref} tag is, a special case for a one-element \code{map} only if the
$\mathit{calls} = \mathit{many}$ tag is.

\section{What an implementation must provide}
\label{sec:eval-brief}

For the experiment to be run, an implementation of the report mode must provide: the dump format
above; the demand-set flag; the hand-over entry assumption as a per-program fact, with the program's
shape read from the write-set analysis; the derivation tags of \S\ref{sec:errors-tags}; one test program
per transfer rule of Table~\ref{tab:transfer} and per summary rule of \S\ref{sec:inf-summarise}, written
to fail before the rule exists; the determinism test's comparison of the dump; a profiling event for
the pass; and the hand-classification sheet with its three columns. No emitted byte changes. The pass
is to be specified as a normative part of the compiler's design before it is built; this thesis is the
argument for that specification, not the specification itself.

\section{Performance, afterwards}
\label{sec:eval-perf}

If the rule is adopted, a second experiment repeats the prototype study's benchmarks
(\S\ref{sec:bg-measure}) with the trie removed and lists written in place---the static variant's
numbers, without its fallback---together with page-size measurements and the direct platform's table
benchmarks and their sweeps. The prototype study gives the expected shape: large gains when building in
a fold and on steady grid ticks, gains hidden by the render walk on TEA table messages, and nothing for
histories, which must copy. That expectation is to be measured, not assumed.
